\abstract{Digital mental-health and cognitive-training studies increasingly collect dense behavioral and physiological data, but improvement is still often summarized by accuracy, mean response time, or symptom scores. This review examines a different class of outcome: the temporal process by which performance or state returns toward a defined reference after a discrete perturbation. We will systematically search biomedical, psychological, engineering, and computing databases and map how studies define the perturbation, reference state, onset of recovery, temporal endpoint, persistence/stability rule, censoring, missingness, within-person normalization, and external validation. The primary objective is a taxonomy of temporal recovery endpoints and their reporting requirements. Because substantial heterogeneity is expected across tasks, populations, and modalities, the synthesis is planned as a systematic scoping review rather than an effect-size meta-analysis. Results will not be used to retroactively justify a particular ARC-ADHD endpoint; instead, the review will identify convergent practices, unresolved ambiguities, and minimum information needed to reproduce a recovery measure.}
